\addtocontents{toc}{\protect\newpage}
\chapter{Joint Distributions, Dependence, and Random Vectors}
\chaptermark{Joint Distributions and Random Vectors}
\label{chap:joint-distributions}

\section{Joint and marginal distributions}
\label{sec:joint-tables}

In the shared-user experiment of Section~\ref{sec:mixtures}, record
the first two responses separately: \(X,Y\in\{0,1\}\), rather than
only their total. Their \emph{joint probability mass function}
\[
  p_{X,Y}(x,y)=\Prob(X=x,Y=y)
\]
records the probability of each pair. With the same equally likely
user probabilities \(0.1\) and \(0.9\), it is the table below. For
example, the two-click cell is
\(p_{X,Y}(1,1)=\tfrac12(0.1)^2+\tfrac12(0.9)^2=0.41\).

\begin{center}
\renewcommand{\arraystretch}{1.3}
\begin{tabular}{@{}lccc@{}}
\toprule
 & \(Y=0\) & \(Y=1\) & Row total\\
\midrule
\(X=0\) & \(0.41\) & \(0.09\) & \(0.50\)\\
\(X=1\) & \(0.09\) & \(0.41\) & \(0.50\)\\
\midrule
Column total & \(0.50\) & \(0.50\) & \(1\)\\
\bottomrule
\end{tabular}
\end{center}

\subsection{Marginalizing a table}

If we only care about the first response, either value of the second is
acceptable. We therefore add across a row:
\begin{equation}
  p_X(x)=\sum_y p_{X,Y}(x,y),\qquad
  p_Y(y)=\sum_x p_{X,Y}(x,y).
  \label{eq:joint-marginals-discrete}
\end{equation}
These are the \emph{marginal} PMFs, written in the margins of the table.
Marginalizing means ignoring a variable by summing over its possible
values, not setting it to zero.

If the two responses instead came from independently selected users,
every interior cell would be \(0.25\). The row and column totals would
still be \(0.5\). Thus the same two marginal distributions can belong
to very different joint distributions.

\subsection{Conditional PMFs and factorization}

Applying conditional probability \eqref{eq:conditional-probability} to
the table gives
\[
  \Prob(Y=1\mid X=1)=\frac{0.41}{0.50}=0.82,\qquad
  \Prob(Y=1\mid X=0)=\frac{0.09}{0.50}=0.18.
\]
For a fixed \(x\), the \emph{conditional PMF} is the selected row
divided by its total. In notation, for \(p_X(x)>0\),
\begin{equation}
  p_{Y\mid X}(y\mid x)=\frac{p_{X,Y}(x,y)}{p_X(x)},
  \qquad
  p_{X,Y}(x,y)=p_X(x)p_{Y\mid X}(y\mid x).
  \label{eq:joint-factorization-discrete}
\end{equation}
The independence condition \eqref{eq:independence}, applied to all
pairs of values, becomes
\begin{equation}
  p_{X,Y}(x,y)=p_X(x)p_Y(y)\quad\text{for every }x,y.
  \label{eq:joint-independence-discrete}
\end{equation}
The table fails this test: \(0.41\ne0.5\cdot0.5\). Once the user's
probability \(P\) is fixed, however, the conditional joint PMF factors.
This is \emph{conditional independence}, written
\(X\independent Y\mid P\). The shared-user effect from Chapter~2
is therefore a distinction between conditional and unconditional
factorization, not a causal effect of one click on the other.

\section{Joint densities: probability over a region}
\label{sec:joint-densities}

For a pair with a joint density, replace the probability table by a density over
the plane. Near \((x,y)\), a small rectangle of side lengths
\(\Delta x,\Delta y\) has probability approximately
\(f_{X,Y}(x,y)\Delta x\Delta y\). Adding small rectangles becomes a
double integral. For a region \(R\),
\[
  \Prob((X,Y)\in R)=\iint_R f_{X,Y}(x,y)\,dx\,dy.
\]
The density is nonnegative, and its integral over the whole plane is one.

\subsection{Two arrival times, put in order}

Suppose two people each independently choose an arrival time uniformly
within a one-hour window. Let \(U,V\) be their times in hours after the
window begins. Their unordered pair is uniform on the unit square:
\(f_{U,V}(u,v)=1\) for \(0<u,v<1\).

Now record the earlier and later times,
\[
  T_1=\min(U,V),\qquad T_2=\max(U,V).
\]
For \(t_1<t_2\), the ordered pair near \((t_1,t_2)\) can arise from
\((U,V)\) near either \((t_1,t_2)\) or \((t_2,t_1)\). These two
equally sized regions contribute equally. Hence
\[
  f_{T_1,T_2}(t_1,t_2)=
  \begin{cases}
    2,&0<t_1<t_2<1,\\
    0,&\text{otherwise}.
  \end{cases}
\]
The allowed triangle has area \(1/2\); multiplying by the density two
gives total probability one. Sorting has introduced dependence: the
later time cannot be smaller than the earlier time.

For instance, both arrivals occur in the first half-hour exactly when
\(T_2\le1/2\). The relevant triangle has area
\(\tfrac12(\tfrac12)^2=1/8\), so
\[
  \Prob(T_2\le1/2)=2\cdot\frac18=\frac14.
\]
This agrees with the probability that both independent original times
fall in the first half of the window.

\subsection{Marginalizing integrates along a slice}

Fix an earlier time \(t\). The later time can range from \(t\) to one,
so the marginal density of the earlier time is
\[
  f_{T_1}(t)=\int_t^1 2\,ds=2(1-t),\qquad 0<t<1.
\]
There are more possible later arrivals when the first arrival is early.
Conversely, if the later arrival is at \(s\), the earlier one can be
anywhere between zero and \(s\):
\[
  f_{T_2}(s)=\int_0^s2\,dt=2s,\qquad 0<s<1.
\]
The general marginalization rules are
\begin{equation}
  f_X(x)=\int_{-\infty}^{\infty}f_{X,Y}(x,y)\,dy,
  \qquad
  f_Y(y)=\int_{-\infty}^{\infty}f_{X,Y}(x,y)\,dx.
  \label{eq:joint-marginals-continuous}
\end{equation}
Although the bounds are written over the whole line, the integrand is
zero outside the joint support. The triangular support is what produced
the variable integration limits in the example.

\begin{figure}[htbp]
  \centering
  \begin{tikzpicture}[x=5.2cm,y=5.2cm]
  \fill[CourseBlue!12] (0,0)--(0,1)--(1,1)--cycle;
  \fill[CourseBlue!45] (0,0)--(0,0.5)--(0.5,0.5)--cycle;
  \draw[CourseBlue,thick] (0,0)--(1,1)--(0,1);
  \draw[->] (-0.02,0)--(1.12,0) node[right]{\(t_1\)};
  \draw[->] (0,-0.02)--(0,1.12) node[above]{\(t_2\)};
  \draw[orange!80!black,very thick] (0.25,0.25)--(0.25,1);
  \draw[orange!80!black,dashed] (0.25,0)--(0.25,0.25);
  \draw[dashed,CourseGray] (0,0.5)--(0.5,0.5);
  \foreach \x/\lab in {0.25/{1/4},0.5/{1/2},1/1}{
    \draw (\x,0)--(\x,-0.018) node[below,font=\small]{\(\lab\)};
  }
  \foreach \y/\lab in {0.5/{1/2},1/1}{
    \draw (0,\y)--(-0.018,\y) node[left,font=\small]{\(\lab\)};
  }
  \node[below left,font=\small] at (0,0){\(0\)};
  \node[font=\small,CourseDark] at (0.54,0.84){Density \(2\)};
  \node[font=\small,CourseGray] at (0.73,0.25){Impossible};
  \node[anchor=west,font=\small,text=orange!80!black] at (1.18,0.80)
    {Fix \(t_1=1/4\)};
  \draw[->,orange!80!black] (1.17,0.74)--(0.28,0.67);
\end{tikzpicture}

  \caption{Every possible ordered pair lies above the diagonal
  \(t_2=t_1\). The dark triangle represents two arrivals before half an
  hour. At a fixed \(t_1=1/4\), the later time ranges along the marked
  vertical segment from \(1/4\) to one. Marginalizing adds density along
  such a segment; conditioning normalizes the segment's density.}
  \label{fig:joint-arrival-region}
\end{figure}

\subsection{A conditional density is a normalized slice}

If the first arrival is at time \(t\), how much longer do we wait for
the second? Dividing the joint density by the earlier-time marginal gives
\[
  f_{T_2\mid T_1}(s\mid t)
  =\frac{2}{2(1-t)}=\frac1{1-t},\qquad t<s<1.
\]
Thus, conditional on the earlier time, the later time is uniform over
the remaining part of the window. At \(t=1/4\), for example,
\[
  \Prob(T_2>3/4\mid T_1=1/4)
  =\int_{3/4}^{1}\frac{1}{1-1/4}\,ds=\frac13.
\]

In general, wherever \(f_X(x)>0\), the conditional-density formula is
\begin{equation}
  f_{Y\mid X}(y\mid x)=\frac{f_{X,Y}(x,y)}{f_X(x)}.
  \label{eq:joint-conditional-density}
\end{equation}
In a smooth density model, conditioning on a shrinking interval around
\(x\) gives this normalized slice. The denominator is a marginal
density, not a point probability. Conditional densities are specified
up to changes on sets of conditioning values with probability zero.

For a pair with a joint density, independence is equivalent to
\(f_{X,Y}(x,y)=f_X(x)f_Y(y)\) almost everywhere. The ordered times fail
this test: both marginals are positive inside \((0,1)\), but their
joint density is zero whenever \(t_1>t_2\).

\subsection{Functions of the pair}
\label{sec:joint-expectations}

Applying the expectation rule \eqref{eq:expectation-function} to the
joint distribution gives
\begin{equation}
  \Expect[g(X,Y)]=
  \begin{cases}
    \displaystyle\sum_x\sum_y g(x,y)p_{X,Y}(x,y),
      &\text{discrete pair},\\[6pt]
    \displaystyle\iint_{\Real^2}g(x,y)f_{X,Y}(x,y)\,dx\,dy,
      &\text{pair with a density}.
  \end{cases}
  \label{eq:joint-function-expectation}
\end{equation}
We require \(\Expect[|g(X,Y)|]<\infty\) for a finite expectation.
For example, the average gap between the ordered arrivals is
\begin{align*}
  \Expect[T_2-T_1]
    &=\int_0^1\int_t^1(s-t)\,2\,ds\,dt\\
    &=\int_0^1(1-t)^2\,dt=\frac13\text{ hour}.
\end{align*}

\section{Correlation and nonlinear dependence}
\label{sec:covariance-correlation}

For variables with finite second moments, expanding the covariance
definition from Section~\ref{sec:variance} gives a convenient calculation
from a joint distribution:
\begin{equation}
  \Cov(X,Y)=\Expect[XY]-\Expect[X]\Expect[Y].
  \label{eq:covariance-product}
\end{equation}
For the click table, \(XY\) is one only in the \((1,1)\) cell, so
\(\Expect[XY]=0.41\) and
\[
  \Cov(X,Y)=0.41-(0.5)(0.5)=0.16.
\]

\subsection{Correlation removes the units}

Converting one measurement from minutes to seconds multiplies its
covariance with another by \(60\), although the relationship has not
changed. To remove this dependence on units, use \emph{correlation}.
For means \(\mu_X,\mu_Y\) and positive standard deviations
\(\sigma_X,\sigma_Y\), it is
\begin{equation}
  \rho_{X,Y}=\frac{\Cov(X,Y)}{\sigma_X\sigma_Y}
  =\Expect\!\left[
      \frac{X-\mu_X}{\sigma_X}\frac{Y-\mu_Y}{\sigma_Y}
    \right].
  \label{eq:correlation}
\end{equation}
This is the average product of the standardized deviations. Changing
minutes to seconds no longer affects the answer. Reversing the sign of
one variable reverses the sign of the correlation. For our clicks,
\(\sigma_X=\sigma_Y=0.5\), hence \(\rho_{X,Y}=0.64\).

To see why correlation lies between \(-1\) and one, put
\(U=(X-\mu_X)/\sigma_X\) and \(V=(Y-\mu_Y)/\sigma_Y\). Both have
mean zero and variance one, so
\[
  0\le\Expect[(U-V)^2]=2-2\rho_{X,Y},\qquad
  0\le\Expect[(U+V)^2]=2+2\rho_{X,Y}.
\]
At \(\rho=1\), the first square has expectation zero, so \(U=V\)
with probability one. At \(\rho=-1\), \(U=-V\) with probability one.
Thus perfect correlation means an exact affine relationship, with
positive or negative slope. Correlation is undefined if either variable
has variance zero: a constant has no fluctuations to standardize.

\subsection{Zero covariance can hide complete dependence}

The distinction between zero covariance and independence noted in
Section~\ref{sec:variance} has a particularly sharp example: one
variable can determine the other completely while their correlation
is zero.

Let \(X\) be uniform on \([-1,1]\), and let \(Y=X^2\). Then
\[
  \Expect[X]=0,\qquad \Expect[Y]=\Expect[X^2]=\frac13,
  \qquad \Expect[XY]=\Expect[X^3]=0.
\]
The odd powers average to zero by symmetry. Thus
\(\Cov(X,Y)=0\), and the correlation is zero since both variances
are positive. Yet knowing \(X\) tells us \(Y\) exactly.

An event calculation also rules out independence. Here
\(Y\le1/4\) is exactly the event \(|X|\le1/2\), so
\[
  \Prob(Y\le1/4\mid |X|\le1/2)=1,
  \qquad \Prob(Y\le1/4)=\frac12.
\]
Correlation detects alignment of signed deviations, not every possible
relationship. An upward-then-downward or U-shaped relationship can
have zero correlation while remaining highly informative.

\section{Conditional means as prediction rules}
\label{sec:conditional-moments}

The conditional moments from Section~\ref{sec:mixtures} also determine
which prediction rule has the smallest mean squared error. Write
\(m(x)=\Expect[Y\mid X=x]\) and \(v(x)=\Var(Y\mid X=x)\), and assume
\(Y\) has a finite second moment. For a prediction \(a(x)\), expand
\(Y-a(x)=(Y-m(x))+(m(x)-a(x))\) inside the group \(X=x\):
\[
  \Expect[(Y-a(x))^2\mid X=x]
  =v(x)+(m(x)-a(x))^2.
\]
The cross term has conditional mean zero. Applying total expectation
\eqref{eq:total-expectation}, for a rule with finite mean squared error,
\[
  \Expect[(Y-a(X))^2]
  =\Expect[v(X)]+\Expect[(m(X)-a(X))^2].
\]
The first term is uncertainty left in \(Y\) after observing \(X\).
The second is the extra error from predicting something other than the
conditional mean. It is nonnegative and vanishes for \(a(X)=m(X)\).
Thus the conditional mean gives the smallest expected squared prediction
error among rules using the same information. In the click table,
\(m(0)=0.18\) and \(m(1)=0.82\), so its error is
\(0.18(0.82)=0.1476\), compared with \(0.25\) for the constant
prediction \(0.5\). This is a property of the probability model;
estimating such a rule from a dataset is a separate statistical problem.

\section{Random vectors and covariance matrices}
\label{sec:random-vectors}

A pair is the simplest example of several measurements made together.
For a single object we might record height, width, temperature, and two
sensor readings. Collect \(d\) measurements into a column vector
\[
  \mathbf X=
  \begin{pmatrix}X_1\\ \vdots\\ X_d\end{pmatrix}\in\Real^d.
\]
This is a \emph{random vector}: all coordinates refer to the same
randomly selected object or experimental repetition. A lowercase
\(\mathbf x\) denotes one realized vector. The joint distribution
specifies probabilities for regions in \(\Real^d\).

The mean vector averages each coordinate:
\[
  \boldsymbol\mu=\Expect[\mathbf X]
  =\begin{pmatrix}\Expect[X_1]\\ \vdots\\ \Expect[X_d]\end{pmatrix}.
\]
To record how every pair of coordinates fluctuates together, arrange
their covariances into a matrix.

\begin{definition}[Covariance matrix]
For a random vector with finite second moments,
\begin{equation}
  \Sigma=\Expect[(\mathbf X-\boldsymbol\mu)
                    (\mathbf X-\boldsymbol\mu)^\top],
  \qquad \Sigma_{ij}=\Cov(X_i,X_j).
  \label{eq:covariance-matrix}
\end{equation}
The superscript \(\top\) denotes transpose: it turns the column of
deviations into a row. Their product is a \(d\times d\) matrix whose
\((i,j)\) entry is the product of the two deviations.
\end{definition}

The diagonal entries are variances; the off-diagonal entries are
covariances. For the two-click vector,
\[
  \boldsymbol\mu=\begin{pmatrix}0.5\\0.5\end{pmatrix},\qquad
  \Sigma=\begin{pmatrix}0.25&0.16\\0.16&0.25\end{pmatrix}.
\]
The matrix is symmetric because \(\Cov(X_i,X_j)=\Cov(X_j,X_i)\).
If the coordinates have positive variances, standardizing each one
gives the correlation matrix, with entries
\(R_{ij}=\Sigma_{ij}/\sqrt{\Sigma_{ii}\Sigma_{jj}}\) and ones on its
diagonal. For the clicks its off-diagonal entries are \(0.64\).

\subsection{A covariance matrix cannot imply a negative variance}

Take a weighted sum \(S=a_1X_1+\cdots+a_dX_d=\mathbf a^\top\mathbf X\).
Expanding its squared centered value gives
\begin{equation}
  \Var(S)=\sum_{i=1}^d\sum_{j=1}^d
            a_i a_j\Cov(X_i,X_j)
         =\mathbf a^\top\Sigma\mathbf a\ge0.
  \label{eq:variance-quadratic-form}
\end{equation}
A symmetric matrix satisfying this inequality for every real vector
\(\mathbf a\) is called \emph{positive semidefinite}. Every covariance
matrix must have this property. It is \emph{positive definite} if the
inequality is strict for every nonzero \(\mathbf a\), meaning that no
nontrivial weighted combination has zero variance.

For example, the symmetric matrix
\[
  M=\begin{pmatrix}1&2\\2&1\end{pmatrix}
\]
cannot be a covariance matrix. It would imply
\(\Var(X_1-X_2)=1+1-2(2)=-2\). Positive diagonal entries alone do not
make a covariance matrix valid. In two dimensions, a symmetric matrix
\[
  \begin{pmatrix}v_1&c\\c&v_2\end{pmatrix}
\]
is positive semidefinite exactly when
\(v_1\ge0\), \(v_2\ge0\), and \(c^2\le v_1v_2\).
For positive variances, the last condition is the correlation bound
\(|\rho|\le1\).

A diagonal covariance matrix means that distinct coordinates are
uncorrelated. It does not generally mean that they are independent:
\((X,X^2)\), with \(X\) uniform on \([-1,1]\), remains a counterexample.

\section{Linear transformations of random vectors}
\label{sec:random-linear-transforms}

Two sensors measuring the same quantity can be summarized by their
average and their difference. Converting units, subtracting a known
bias, or combining features into a score are other examples of linear
combinations. We want rules for the means and covariances of the results.

Let \(A\) be a fixed \(k\times d\) matrix and \(\mathbf b\in\Real^k\)
a fixed vector, and set
\[
  \mathbf W=A\mathbf X+\mathbf b.
\]
Row \(i\) of \(A\) supplies the weights for the \(i\)-th new variable.
Linearity of expectation gives
\begin{equation}
  \Expect[\mathbf W]=A\boldsymbol\mu+\mathbf b.
  \label{eq:linear-mean-vector}
\end{equation}
Subtracting this mean yields
\(\mathbf W-\Expect[\mathbf W]=A(\mathbf X-\boldsymbol\mu)\). Hence
\begin{align}
  \Cov(\mathbf W)
    &=\Expect[A(\mathbf X-\boldsymbol\mu)
                  (\mathbf X-\boldsymbol\mu)^\top A^\top]\notag\\
    &=A\Sigma A^\top.
  \label{eq:linear-covariance-vector}
\end{align}
Here \(\Cov(\mathbf W)\) denotes the covariance matrix of a vector.
These rules need finite second moments, not independence or Gaussian
distributions.

\begin{example}[Average error and disagreement between sensors]
Let \(X_1,X_2\) be centered sensor errors with variances one and
covariance \(0.8\). Define
\[
  M=\frac{X_1+X_2}{2},\qquad D=X_1-X_2,\qquad
  A=\begin{pmatrix}1/2&1/2\\1&-1\end{pmatrix}.
\]
Both new means are zero, and matrix multiplication gives
\[
  \Cov\!\begin{pmatrix}M\\D\end{pmatrix}
  =A\begin{pmatrix}1&0.8\\0.8&1\end{pmatrix}A^\top
  =\begin{pmatrix}0.9&0\\0&0.4\end{pmatrix}.
\]
Positively correlated errors often move together, so averaging them
removes little of the shared error. Subtracting them cancels that
shared variation. With independent unit-variance errors, the two
variances would instead be \(0.5\) and \(2\).
\end{example}

The zero covariance between \(M\) and \(D\) follows from the equal
original variances. It does not by itself establish independence. The
Gaussian model below supplies an important case where it does.

\section{The multivariate Gaussian}
\label{sec:multivariate-gaussian}

\subsection{Construct correlated readings from independent noise}

Let \(Z_1,Z_2\) be independent standard Gaussian variables. Consider
the sensor-error model
\[
  X_1=Z_1,\qquad X_2=0.8Z_1+0.6Z_2.
\]
Both sensors contain the disturbance \(Z_1\); the second also has an
independent disturbance \(Z_2\). The known rules for independent
Gaussian sums give
\[
  \Expect[X_1]=\Expect[X_2]=0,\qquad
  \Var(X_1)=1,\qquad \Var(X_2)=0.8^2+0.6^2=1,
\]
and
\[
  \Cov(X_1,X_2)=0.8\Var(Z_1)+0.6\Cov(Z_1,Z_2)=0.8.
\]
Thus both marginal distributions are standard Gaussian, but their
joint distribution is correlated. It is a \emph{bivariate Gaussian}.

The construction extends to any number of measurements.

\begin{definition}[Multivariate Gaussian]
Let \(\mathbf Z\) have \(m\) independent standard Gaussian coordinates,
and let \(A\) be a fixed \(d\times m\) matrix. A vector of the form
\[
  \mathbf X=\boldsymbol\mu+A\mathbf Z
\]
has a multivariate Gaussian distribution, written
\[
  \mathbf X\sim\mathcal N_d(\boldsymbol\mu,\Sigma),
  \qquad \Sigma=AA^\top.
\]
Its mean vector is \(\boldsymbol\mu\), and its covariance matrix is
\(\Sigma\).
\end{definition}

The independent noises supply the randomness; the rows of \(A\) specify
which noises each measurement shares. Equivalently, a vector is jointly
Gaussian if every real linear combination of its coordinates is a
univariate Gaussian, allowing a constant when the variance is zero.
For the construction above this follows because each combination is
itself a sum of independent Gaussian noises. In particular,
\begin{equation}
  \mathbf a^\top\mathbf X
  \sim\mathcal N(\mathbf a^\top\boldsymbol\mu,
                 \mathbf a^\top\Sigma\mathbf a).
  \label{eq:gaussian-linear-combination}
\end{equation}

\subsection{From circular to elliptical density contours}

Independence gives the joint density of two standard noises:
\[
  f_{Z_1,Z_2}(z_1,z_2)
  =\phi(z_1)\phi(z_2)
  =\frac1{2\pi}\exp\!\left[-\frac{z_1^2+z_2^2}{2}\right].
\]
It has the same value wherever \(z_1^2+z_2^2\) is constant, so its
equal-density contours are circles centered at zero. These contours
describe density values, not a claim that all observations lie inside
a particular circle.

For a correlation \(-1<\rho<1\), replace the numerical sensor model by
\[
  X_1=Z_1,\qquad X_2=\rho Z_1+\sqrt{1-\rho^2}\,Z_2.
\]
The inverse relations are
\(z_1=x_1\) and
\(z_2=(x_2-\rho x_1)/\sqrt{1-\rho^2}\).
The transformation scales areas by \(\sqrt{1-\rho^2}\), so preserving
probabilities requires dividing the density by this factor. Substitution
then gives
\begin{equation}
  f_{X_1,X_2}(x_1,x_2)
  =\frac1{2\pi\sqrt{1-\rho^2}}
    \exp\!\left[-\frac{x_1^2-2\rho x_1x_2+x_2^2}
                         {2(1-\rho^2)}\right].
  \label{eq:bivariate-standard-density}
\end{equation}
The cross term changes which pairs are likely. With positive \(\rho\),
equal-sign deviations are favored; with negative \(\rho\), opposite-sign
deviations are favored. Both marginals are still \(\mathcal N(0,1)\).

\begin{figure}[htbp]
  \centering
  \begin{tikzpicture}
\foreach \r/\offset/\heading in {0/0/{\rho=0},0.8/4.9/{\rho=0.8},-0.8/9.8/{\rho=-0.8}}{
  \begin{axis}[
    at={(\offset cm,0)},anchor=south west,
    width=4.8cm,height=4.8cm,scale only axis=false,
    axis equal image,axis lines=middle,
    xmin=-2.6,xmax=2.6,ymin=-2.6,ymax=2.6,
    xtick={-2,0,2},ytick={-2,0,2},
    xlabel={\(x_1\)},ylabel={\(x_2\)},
    title={\(\heading\)},title style={font=\small},
    tick label style={font=\scriptsize},label style={font=\small},
    grid=major,grid style={black!6},
    domain=0:360,samples=121,
  ]
    \addplot[CourseBlue,thick] ({cos(x)},{\r*cos(x)+sqrt(1-\r*\r)*sin(x)});
    \addplot[CourseBlue,thick] ({2*cos(x)},{2*(\r*cos(x)+sqrt(1-\r*\r)*sin(x))});
  \end{axis}
}
\end{tikzpicture}

  \caption{Equal-density contours for three bivariate Gaussian models.
  In every panel both individual features have mean zero and variance
  one. Correlation changes the joint shape, which separate histograms
  of the features cannot reveal. Each curve joins points with the same
  density; the outer curve has lower density than the inner one.}
  \label{fig:gaussian-correlation-contours}
\end{figure}

The ellipse directions can be understood without a matrix decomposition.
For \(\rho=0.8\), use the rotated coordinates
\[
  U=\frac{X_1+X_2}{\sqrt2},\qquad
  V=\frac{X_1-X_2}{\sqrt2}.
\]
Their variances are \(1+\rho=1.8\) and \(1-\rho=0.2\), respectively,
and their covariance is zero. In these coordinates the expression in
the exponent becomes
\[
  q=\frac{u^2}{1.8}+\frac{v^2}{0.2}.
\]
For fixed \(q=c>0\), the ellipse extends by \(\sqrt{1.8c}\) along
the equal-sign direction and \(\sqrt{0.2c}\) along the opposite-sign
direction. The first is three times longer than the second. Shared
sensor fluctuations stretch the distribution along the diagonal.

\subsection{The density in matrix notation}

For a positive-definite covariance matrix, the general Gaussian density
is
\begin{equation}
  f_{\mathbf X}(\mathbf x)
  =\frac{1}{(2\pi)^{d/2}\sqrt{\det\Sigma}}
    \exp\!\left[-\frac12(\mathbf x-\boldsymbol\mu)^\top
      \Sigma^{-1}(\mathbf x-\boldsymbol\mu)\right].
  \label{eq:multivariate-gaussian-density}
\end{equation}
Here \(\Sigma^{-1}\) is the inverse matrix, and \(\det\Sigma\) is its
determinant. These are the matrix versions of dividing by variance and
correcting the density for a change of scale. For an invertible square
\(A\), the transformation \(\mathbf x=\boldsymbol\mu+A\mathbf z\)
scales volumes by \(|\det A|=\sqrt{\det\Sigma}\), while
\[
  \sum_i z_i^2
  =(\mathbf x-\boldsymbol\mu)^\top\Sigma^{-1}
      (\mathbf x-\boldsymbol\mu).
\]
This explains both parts of the density formula from the independent
standard-noise density.

In two dimensions,
\[
  \Sigma=\begin{pmatrix}
    \sigma_1^2&\rho\sigma_1\sigma_2\\
    \rho\sigma_1\sigma_2&\sigma_2^2
  \end{pmatrix},\qquad
  \det\Sigma=\sigma_1^2\sigma_2^2(1-\rho^2).
\]
For positive \(\sigma_1,\sigma_2\) and \(|\rho|<1\), putting
\(u=(x_1-\mu_1)/\sigma_1\), \(v=(x_2-\mu_2)/\sigma_2\) gives
\[
  (\mathbf x-\boldsymbol\mu)^\top\Sigma^{-1}
      (\mathbf x-\boldsymbol\mu)
  =\frac{u^2-2\rho uv+v^2}{1-\rho^2}.
\]
For a diagonal covariance matrix this reduces to a sum of squared
standardized deviations, and the density factors into the univariate
Gaussian densities.

The quadratic expression is called the \emph{squared Mahalanobis
distance}. It measures how unusual a displacement is relative to the
model's spread and dependence. In the unit-variance model with
\(\rho=0.8\), the points \((1,1)\) and \((1,-1)\) have the same
ordinary distance from the origin, but
\[
  q(1,1)=\frac{2-1.6}{0.36}=\frac{10}{9},\qquad
  q(1,-1)=\frac{2+1.6}{0.36}=10.
\]
Disagreement between the sensors is much less typical than a shared
positive error. Euclidean distance alone misses that distinction.

\subsection{Marginals, transformations, and the independence exception}

Selecting some coordinates is a linear transformation. Their marginal
distribution is Gaussian with the corresponding entries of the mean
vector and the corresponding rows and columns of the covariance matrix.
In particular, \(X_i\sim\mathcal N(\mu_i,\Sigma_{ii})\).

More generally, for fixed \(B\) and \(\mathbf b\),
\begin{equation}
  B\mathbf X+\mathbf b
  \sim\mathcal N(B\boldsymbol\mu+\mathbf b, B\Sigma B^\top).
  \label{eq:gaussian-vector-transform}
\end{equation}
The result is Gaussian because substituting
\(\mathbf X=\boldsymbol\mu+A\mathbf Z\) gives another linear
transformation of independent standard Gaussian noises. The moment
rules in Section~\ref{sec:random-linear-transforms} now describe the
\emph{whole distribution}, not just two summaries of it.

For jointly Gaussian variables, zero covariance \emph{does} imply
independence. In the nondegenerate case, a diagonal covariance matrix
makes the joint density a product of marginal densities. Therefore the
average and difference in the sensor example are independent under
the joint Gaussian assumption, as well as uncorrelated.

It is not enough for each marginal to be Gaussian. To see why, take
\(X\sim\mathcal N(0,1)\) and an independent random sign \(S\), equal to
\(-1\) or one with probability \(1/2\). Set \(Y=SX\). Symmetry makes
\(Y\sim\mathcal N(0,1)\), and
\[
  \Cov(X,Y)=\Expect[SX^2]=\Expect[S]\Expect[X^2]=0.
\]
Yet \(|Y|=|X|\) always, so they are dependent. The pair is not jointly
Gaussian: \(X+Y\) is exactly zero with probability \(1/2\), but is not
a constant. A nondegenerate Gaussian has no such point mass.

Finally, a Gaussian covariance matrix can be singular. If
\(X_2=X_1\sim\mathcal N(0,1)\), then
\[
  \Sigma=\begin{pmatrix}1&1\\1&1\end{pmatrix}.
\]
All pairs lie on a line. This is a valid Gaussian random vector, but it
has no density with respect to two-dimensional area. Formula
\eqref{eq:multivariate-gaussian-density} requires a positive-definite
matrix; it cannot be used when the determinant is zero.

\section{What one Gaussian measurement tells us about another}
\label{sec:conditional-gaussian}

Return to \(X_1=Z_1\), \(X_2=0.8Z_1+0.6Z_2\). Before seeing the
first sensor, the second error has variance one. After observing
\(X_1=x\), the shared noise \(Z_1\) is known and only the independent
noise \(Z_2\) remains:
\[
  X_2\mid X_1=x\sim\mathcal N(0.8x,0.36).
\]
If the first sensor reads two units too high, the conditional mean
error of the second is \(1.6\), with standard deviation \(0.6\).
Marginalizing over the first sensor and conditioning on its reading are
different operations: the marginal distribution is still
\(\mathcal N(0,1)\).

For the general bivariate Gaussian with positive marginal variances and
\(|\rho|<1\), write
\[
  X=\mu_X+\sigma_X Z_1,\qquad
  Y=\mu_Y+\sigma_Y(\rho Z_1+\sqrt{1-\rho^2}Z_2).
\]
Observing \(X=x\) fixes \(Z_1=(x-\mu_X)/\sigma_X\). Therefore
\begin{equation}
  Y\mid X=x\sim\mathcal N\!\left(
    \mu_Y+\rho\frac{\sigma_Y}{\sigma_X}(x-\mu_X),
    \sigma_Y^2(1-\rho^2)\right).
  \label{eq:conditional-bivariate-gaussian}
\end{equation}
Equivalently, the conditional mean's slope is
\(\Cov(X,Y)/\Var(X)\). Positive correlation shifts the conditional
mean in the direction of the observation; negative correlation shifts
it in the opposite direction. The conditional variance is constant in
\(x\) for this Gaussian model. General conditional distributions need
not have a linear mean or a constant variance.

\begin{figure}[H]
  \centering
  \begin{tikzpicture}
\begin{axis}[
  width=0.94\textwidth,height=6.4cm,
  axis lines=left,xmin=-3.5,xmax=4,ymin=0,ymax=0.86,
  xlabel={Second sensor error \(x_2\)},ylabel={Density},
  tick label style={font=\small},label style={font=\small},
  legend style={draw=none,font=\small,at={(0.5,0.99)},anchor=north},
  grid=major,grid style={black!8},domain=-3.5:4,samples=180,
]
  \addplot[CourseGray,thick,dashed] {exp(-x^2/2)/sqrt(2*pi)};
  \addlegendentry{Marginal: \(\mathcal N(0,1)\)}
  \addplot[CourseBlue,very thick] {exp(-(x-1.6)^2/(2*0.36))/(0.6*sqrt(2*pi))};
  \addlegendentry{Given \(X_1=2\): \(\mathcal N(1.6,0.36)\)}
\end{axis}
\end{tikzpicture}

  \caption{Before observing the first sensor, the second error has
  distribution \(\mathcal N(0,1)\). After observing \(X_1=2\), it has
  distribution \(\mathcal N(1.6,0.36)\). Knowledge of the shared
  disturbance shifts the mean and reduces the remaining spread.}
  \label{fig:gaussian-conditional-density}
\end{figure}

As a check using total variance \eqref{eq:total-variance}, the numerical
model gives \(1=0.36+0.64\): residual variance plus variance of the
conditional mean. If \(\rho=0\), observing
\(X\) changes neither the mean nor the variance, consistent with
Gaussian independence. As \(|\rho|\) approaches one, the remaining
variance approaches zero. At \(|\rho|=1\), the relationship is exact
and the conditional distribution is a point mass.

\section{Gaussian class-conditional models}
\label{sec:gaussian-classification}

Let \(C\) denote the class label, taking values \(1,\ldots,K\). Use
Gaussian class-conditional densities in \eqref{eq:bayesian-classifier}:
\[
  \Prob(C=k)=\pi_k,\qquad
  \mathbf X\mid C=k\sim\mathcal N_d(\boldsymbol\mu_k,\Sigma_k).
\]
For positive-definite \(\Sigma_k\), let \(f_k\) be the corresponding
Gaussian density. The new modeling choice is the covariance structure
\emph{within each class}.

\subsection{The relationship between features can identify the class}

Consider two equally likely operating modes of a device. In mode
\(C=+\), its two sensor errors tend to move together; in mode \(C=-\),
they tend to move in opposite directions. Use
\[
  \boldsymbol\mu_+=\boldsymbol\mu_-=
  \begin{pmatrix}0\\0\end{pmatrix},\qquad
  \Sigma_+=\begin{pmatrix}1&0.8\\0.8&1\end{pmatrix},\qquad
  \Sigma_-=\begin{pmatrix}1&-0.8\\-0.8&1\end{pmatrix}.
\]
Each feature alone has distribution \(\mathcal N(0,1)\) in both
classes. Its value alone therefore leaves the class probabilities at
\(1/2\). The distinguishing information is entirely in the pair.

The determinants are both \(0.36\), so the density normalization
constants cancel in the likelihood ratio. Subtracting the exponents
in \eqref{eq:bivariate-standard-density} gives
\[
  \log\frac{f_+(x_1,x_2)}{f_-(x_1,x_2)}
  =\frac{2(0.8)}{1-0.8^2}x_1x_2=\frac{40}{9}x_1x_2.
\]
With equal prior probabilities, mode \(+\) is more probable when
\(x_1x_2>0\), and mode \(-\) is more probable when \(x_1x_2<0\).
At \((1,1)\),
\[
  \Prob(C=+\mid\mathbf X=(1,1))
  =\frac{1}{1+e^{-40/9}}\approx0.9884.
\]
At \((1,-1)\), this probability is approximately \(0.0116\).
The individual feature values have equal marginal densities in both
classes; their agreement or disagreement changes the conclusion.

If we incorrectly replace each class density by the product of its
feature marginals, both models become \(\phi(x_1)\phi(x_2)\). Then
every pair receives posterior class probability \(1/2\).
Assuming conditional independence has erased all the useful information
in this example.

\subsection{What the Gaussian decision rule uses}

For positive class priors, maximizing the posterior is equivalent to
maximizing the log-score
\[
  s_k(\mathbf x)=\log\pi_k-\frac12\log\det\Sigma_k
   -\frac12(\mathbf x-\boldsymbol\mu_k)^\top
           \Sigma_k^{-1}(\mathbf x-\boldsymbol\mu_k).
\]
The omitted term \(-\tfrac d2\log(2\pi)\) is the same for every class.
The three remaining terms reflect the class's prevalence, the volume
over which its density is spread, and the observation's displacement
relative to the class's covariance. A broad class cannot receive the
same peak density as a narrow class with the same total probability.

If all classes share a common covariance \(\Sigma\), expanding the
quadratic expression shows that the term
\(-\tfrac12\mathbf x^\top\Sigma^{-1}\mathbf x\) is also common.
The class-dependent part becomes
\[
  \mathbf x^\top\Sigma^{-1}\boldsymbol\mu_k
  -\frac12\boldsymbol\mu_k^\top\Sigma^{-1}\boldsymbol\mu_k
  +\log\pi_k,
\]
which is linear in the observed features, plus a constant. Comparing
two classes therefore gives a linear boundary. If the covariance
matrices differ, quadratic terms generally remain, as they did in the
two-mode example.

These conclusions come from specifying a joint probability model and
conditioning on an observation. How to learn the means, covariance
matrices, and class probabilities from data belongs to statistical
estimation. Here the important point is that a feature's usefulness
can depend on its relationship with other features, not just its own
distribution.
